% ============================================================================
% FinABSA--CafeF empirical validation section
%
% This is a section fragment, not a standalone LaTeX document. The main
% document should load at least:
%   \usepackage{amsmath,amssymb,booktabs,graphicx,longtable,array,ragged2e}
%   \usepackage{float}
% The figure files referenced below should be placed in images/.
% All numerical results are from the strict 208-row CafeF run.
% ============================================================================

\section{Empirical Validation and Statistical Testing}
\label{sec:validation}

This section evaluates whether sentiment predictions generated by FinABSA are associated with subsequent stock-price movements in a one-month Vietnamese financial-news sample. The analysis is deliberately separated into two layers. The first layer validates the data contract and constructs a reproducible article--ticker event sample. The second layer performs downstream financial tests, including event-study abnormal returns, sentiment-group contrasts, a news-day control comparison, a two-way fixed-effects regression, and lagged sensitivity analyses. Unless explicitly stated otherwise, all reported results use the strict entity-linking sample and the market model estimated with the VNINDEX as the market proxy.

The empirical objective is not to claim that news causes returns. Rather, it is to test whether model-generated sentiment contains an observable association with abnormal price movements after entity linking, time alignment, and market adjustment. This distinction is important because the model output is observational, the sample covers only October 2022, and both the news-generation process and the entity-linking process may be correlated with market conditions.

\subsection{Data Construction and Coverage}
\label{subsec:validation-data}

\paragraph{CafeF corpus.}
We collected and parsed 1,472 CafeF articles published between 1 and 31 October 2022. The crawler used the October sitemap partitions, retained finance-relevant pages, and stored the article URL, title, summary, body, author, publication timestamp, update timestamp, and parsing status. Each article was assigned a stable identifier derived from its URL. The corpus is therefore a fixed historical snapshot rather than a continuously updated news feed.

A single article can mention several companies. Because the ABSA model predicts sentiment toward a specified target rather than discovering all targets itself, the unit passed to the model is an article--ticker pair. The broad candidate mapping produced 512 such pairs. The primary analysis uses 208 pairs whose entity-linking confidence was at least 0.95. These 208 observations correspond to 50 distinct security tickers in the prediction file; the subsequent market panel contains 51 tickers after the available price coverage and event construction are applied.

\begin{table}[htbp]
\centering
\caption{Observed sample and prediction-contract statistics.}
\label{tab:validation_sample}
\begin{tabular}{lrr}
\toprule
\textbf{Quantity} & \textbf{Count} & \textbf{Definition} \\
\midrule
CafeF articles & 1,472 & Unique parsed articles in October 2022 \\
Broad article--ticker pairs & 512 & All automatically detected candidates \\
Strict article--ticker pairs & 208 & Entity confidence $\geq 0.95$ \\
Predictions aligned to strict input & 208 & One prediction per strict pair \\
Neutral predictions & 180 & Hard label ``neutral'' \\
Positive predictions & 18 & Hard label ``positive'' \\
Negative predictions & 10 & Hard label ``negative'' \\
Usable crisis-window event rows & 200 & Rows with matched event-period return \\
Ticker--day panel observations & 152 & One market outcome per ticker--day \\
\bottomrule
\end{tabular}
\end{table}

\paragraph{Entity linking and target masking.}
The broad file preserves every candidate detected by the entity-linking rules. The strict file retains the subset with confidence at least 0.95. In the observed data, 192 strict rows were obtained through a curated company-alias match with confidence 0.98, and 16 rows through a ticker-context rule with confidence 0.95. The remaining 304 broad-only rows were explicit uppercase strings assigned confidence 0.90 and were excluded from the primary analysis because uppercase abbreviations can refer to non-company terms, indices, or unrelated financial vocabulary.

For each target, the target surface form is replaced by \texttt{Target}, while other detected entities are replaced by \texttt{Other}. This creates one target-specific input per article--ticker pair and prevents the model from using the target's surface name as a shortcut. For example, a headline mentioning both EIB and CTG generates two distinct masked inputs, one for each target. The prediction file supplied for this analysis contains a hard classification output for every strict input and is therefore sufficient for label-based aggregation; it does not contain calibrated class probabilities.

\begin{figure}[htbp]
\centering
\includegraphics[width=0.78\textwidth]{images/prediction_distribution.png}
\caption{Distribution of FinABSA predictions in the strict CafeF sample. The sample is strongly imbalanced toward the neutral class: 180 of 208 predictions are neutral, compared with 18 positive and 10 negative predictions. The imbalance is considered explicitly when interpreting group-level estimates and statistical power.}
\label{fig:prediction-distribution}
\end{figure}

\paragraph{Market data.}
Daily OHLCV observations were retrieved through \texttt{vnstock} for the identified securities and two reference series, VNINDEX and VN30. The stored price file contains 9,182 rows for 55 unique symbols, including the two reference series; 53 symbols are individual securities. The observed date range is 4 May--30 December 2022. Most symbols have 171 daily observations, while the minimum coverage is 60 observations. The estimation sample therefore spans substantially more than the October event month and permits pre-event market-model estimation.

The primary market proxy is VNINDEX. Log returns are computed as
\[
 r_{i,t}=100\times\left[\log(P_{i,t})-\log(P_{i,t-1})\right],
\]
where $P_{i,t}$ denotes the recorded closing price. The analysis uses the price series as stored in the repository. Because the source file does not provide a separate verified corporate-action audit trail, the interpretation of the magnitude of returns is conditional on the adjustment and coverage properties of the downloaded series.

\begin{table}[htbp]
\centering
\caption{Market-data coverage used in the event study.}
\label{tab:market-coverage}
\begin{tabular}{lr}
\toprule
\textbf{Metric} & \textbf{Value} \\
\midrule
Unique symbols in price file & 55 \\
Reference symbols & VNINDEX, VN30 \\
Security symbols & 53 \\
Total OHLCV rows & 9,182 \\
First available date & 4 May 2022 \\
Last available date & 30 December 2022 \\
Median observations per symbol & 171 \\
Minimum observations per symbol & 60 \\
Maximum observations per symbol & 171 \\
Symbols with fitted market model & 54 \\
\bottomrule
\end{tabular}
\end{table}

\subsection{Prediction Validation and Aggregation}
\label{subsec:prediction-validation}

Before any financial calculation, the prediction file is validated against the strict input contract. Alignment is performed by \texttt{sample\_id}, not by row order. The validation stage requires one prediction for every input identifier, rejects duplicate identifiers, records unexpected identifiers, and canonicalizes the model output to one of three labels: positive, neutral, or negative. The final run has 208 input rows and 208 prediction rows, with no missing predictions, no unexpected predictions, no duplicate input IDs, no duplicate prediction IDs, and no unknown labels.

Because the submitted model output contains hard labels rather than class probabilities, the primary sentiment score is defined as
\[
 S_{i,t}=\begin{cases}
 +1, & \text{if the prediction is positive},\\
 0, & \text{if the prediction is neutral},\\
 -1, & \text{if the prediction is negative}.
 \end{cases}
\]
This is a label-based score, not a softmax probability difference. For each ticker--day, the pipeline also records the mean score, label shares, article count, and the number of article--ticker rows. The article count is transformed as $\log(1+\text{article\_count})$ in the regression.

The aggregation step is necessary because several articles may refer to the same ticker on the same date, whereas the market supplies only one return observation for that ticker--day. Repeating the same return for every article would artificially increase the effective sample size. Accordingly, descriptive article-level event summaries and the ticker--day panel are reported separately. The panel uses one observation per ticker--day and therefore contains 152 observations.

\subsection{Time Alignment}
\label{subsec:time-alignment}

The data contain daily prices rather than intraday transactions. The implemented alignment rule therefore maps each article's normalized publication date to the first trading date strictly after that calendar date. In other words, an article published on a trading day is evaluated against the next available trading session, and an article published on a weekend or holiday is evaluated against the next trading session. This rule avoids attributing a same-day close to information whose exact intraday dissemination and market absorption cannot be observed in the daily data.

This implementation differs from a fixed 15:00 market-close rule: no intraday cutoff is imposed in the completed run. Timestamp fields are retained for traceability, but the event mapping is performed at daily frequency. Eight of the 208 aligned prediction rows do not contribute to the crisis-window event summary because a usable matching stock return is unavailable after the ticker/date join; the resulting crisis event table contains 200 usable rows.

\subsection{Event-Study Design}
\label{subsec:event-study}

We use a standard market model following the event-study framework of Brown and Warner~\cite{brownwarner1985}:
\[
 R_{i,t}=\alpha_i+\beta_iR_{m,t}+\varepsilon_{i,t},
\]
where $R_{i,t}$ is the log return of security $i$, $R_{m,t}$ is the VNINDEX return, and $(\hat{\alpha}_i,\hat{\beta}_i)$ are estimated using observations before 1 October 2022. The expected return is
\[
 \widehat{E}(R_{i,t})=\hat{\alpha}_i+\hat{\beta}_iR_{m,t},
\]
and abnormal return is
\[
 AR_{i,t}=R_{i,t}-\widehat{E}(R_{i,t}).
\]
The market model is estimated separately for each symbol. At least 20 pre-event observations are required for a fitted market model; symbols with less coverage use the repository's documented fallback parameters. In the final run, 54 symbols have market-model records, and the fitted models are used to calculate event-period abnormal returns.

For each usable ticker--day event, cumulative abnormal returns are computed over three windows:
\[
 CAR_i(T_1,T_2)=\sum_{\tau=T_1}^{T_2}AR_{i,t+\tau},
\]
with windows $[0,0]$, $[0,+1]$, and $[0,+3]$. The CAR tests use the 152 ticker--day observations for which the corresponding event-window calculation is available. For each window and sentiment group, the pipeline reports the sample mean, median, one-sample $t$-test against zero, a two-sided sign-flip permutation test, and a nonparametric bootstrap 95\% confidence interval. Both permutation and bootstrap calculations use 5,000 replications and fixed random seeds for reproducibility.

\begin{figure}[htbp]
\centering
\includegraphics[width=0.82\textwidth]{images/abnormal_return_by_label.png}
\caption{Mean same-day abnormal return by predicted sentiment. The positive group has the largest point estimate, while the negative group has a negative point estimate. However, the group sizes are small and the inferential tests and confidence intervals do not establish statistically reliable differences.}
\label{fig:ar-by-label}
\end{figure}

\begin{table}[htbp]
\centering
\scriptsize
\caption{Event-study CAR tests by sentiment group and event window. The reported $p$-values are two-sided. Bootstrap intervals are percentile 95\% confidence intervals.}
\label{tab:car-tests}
\begin{tabular}{llrrrrrrrr}
\toprule
\textbf{Window} & \textbf{Group} & \textbf{N} & \textbf{Mean CAR} & \textbf{Median} & \textbf{$t$} & \textbf{$p_t$} & \textbf{$p_{perm}$} & \textbf{CI low} & \textbf{CI high} \\
\midrule
$[0,0]$ & All & 152 & 0.0990 & -0.0337 & 0.3824 & 0.7027 & 0.6937 & -0.3944 & 0.6010 \\
$[0,0]$ & Negative & 10 & -0.3665 & -1.2829 & -0.3542 & 0.7314 & 0.7351 & -2.2612 & 1.5844 \\
$[0,0]$ & Neutral & 127 & 0.0402 & -0.0655 & 0.1416 & 0.8876 & 0.8902 & -0.5153 & 0.5901 \\
$[0,0]$ & Positive & 15 & 0.9068 & 0.4796 & 1.1224 & 0.2806 & 0.2835 & -0.5604 & 2.4571 \\
\midrule
$[0,+1]$ & All & 152 & -0.0369 & 0.0918 & -0.1064 & 0.9154 & 0.9152 & -0.7098 & 0.6092 \\
$[0,+1]$ & Negative & 10 & 0.1044 & -0.4044 & 0.1019 & 0.9211 & 0.9098 & -1.7411 & 2.0379 \\
$[0,+1]$ & Neutral & 127 & 0.1073 & 0.2594 & 0.2725 & 0.7857 & 0.7886 & -0.6550 & 0.8759 \\
$[0,+1]$ & Positive & 15 & -1.3522 & -1.6099 & -1.5819 & 0.1360 & 0.1314 & -2.9966 & 0.3553 \\
\midrule
$[0,+3]$ & All & 152 & -0.2649 & -0.6823 & -0.5407 & 0.5895 & 0.5971 & -1.2305 & 0.6768 \\
$[0,+3]$ & Negative & 10 & -2.6645 & -2.8650 & -1.2247 & 0.2518 & 0.2492 & -6.6793 & 1.2969 \\
$[0,+3]$ & Neutral & 127 & 0.1476 & 0.0022 & 0.2741 & 0.7845 & 0.7784 & -0.8815 & 1.2090 \\
$[0,+3]$ & Positive & 15 & -2.1575 & -1.9846 & -1.7703 & 0.0984 & 0.0978 & -4.4665 & 0.1626 \\
\bottomrule
\end{tabular}
\end{table}

\begin{figure}[htbp]
\centering
\includegraphics[width=0.82\textwidth]{images/car_by_window.png}
\caption{Mean cumulative abnormal returns across the three event windows. Error bars represent bootstrap 95\% confidence intervals for the all-event sample. The intervals cross zero for all windows, indicating that the point estimates are not statistically distinguishable from zero in this sample.}
\label{fig:car-by-window}
\end{figure}

\paragraph{Event-study assessment.}
The same-day estimate for all events is 0.0990 percentage points, with a one-sample $t$-test $p$-value of 0.7027 and a sign-flip permutation $p$-value of 0.6937. The bootstrap 95\% interval, $[-0.3944,0.6010]$, also contains zero. The positive group has a descriptive same-day mean of 0.9068 percentage points, whereas the negative group has a descriptive mean of -0.3665 percentage points; neither estimate is statistically significant, and both intervals are wide. Over longer windows, the all-event mean becomes slightly negative, but the $[0,+1]$ and $[0,+3]$ tests likewise provide no evidence against a zero mean.

The positive-group estimate changes sign after the event day, moving from $+0.9068$ on $[0,0]$ to $-1.3522$ on $[0,+1]$ and $-2.1575$ on $[0,+3]$. This instability is not consistent with a simple monotonic price-response interpretation. It is also vulnerable to the positive group's small sample size ($N=15$ at the ticker--day level). The correct conclusion is therefore that the event-study results are suggestive descriptive patterns rather than evidence of a stable sentiment effect.

\subsection{News-Day Versus No-News-Day Control}
\label{subsec:news-control}

The event study is complemented by a control comparison that asks whether eligible tickers exhibit different abnormal returns on days with an identified CafeF event relative to eligible ticker--days with no identified event. The control sample is restricted to the same October crisis window and to tickers that appear in the event sample. There are 152 news ticker--days and 934 eligible no-news ticker--days.

We compare the two distributions using a Welch unequal-variance test, a pooled-label permutation test with 5,000 permutations, and a bootstrap confidence interval for the difference in means. The estimand is
\[
 \Delta=\overline{AR}_{\text{news}}-\overline{AR}_{\text{no-news}}.
\]

\begin{table}[htbp]
\centering
\caption{News-day versus eligible no-news-day abnormal returns.}
\label{tab:news-control}
\begin{tabular}{lrrrrrrrr}
\toprule
\textbf{Comparison} & \textbf{$N_N$} & \textbf{$N_0$} & \textbf{News mean} & \textbf{No-news mean} & \textbf{Difference} & \textbf{Welch $p$} & \textbf{Perm. $p$} & \textbf{Bootstrap 95\% CI} \\
\midrule
News $-$ no-news & 152 & 934 & 0.0990 & -0.3509 & 0.4499 & 0.1074 & 0.0988 & $[-0.0901,1.0040]$ \\
\bottomrule
\end{tabular}
\end{table}

\begin{figure}[htbp]
\centering
\includegraphics[width=0.78\textwidth]{images/news_day_control.png}
\caption{Difference in mean abnormal return between news ticker--days and eligible no-news ticker--days. The point estimate is positive, but the bootstrap interval crosses zero.}
\label{fig:news-day-control}
\end{figure}

The point estimate indicates that news days have, on average, 0.4499 percentage points higher abnormal returns than eligible no-news days. Nevertheless, the Welch test gives $p=0.1074$, the permutation test gives $p=0.0988$, and the bootstrap interval includes zero. Thus, the control comparison does not provide conventional 5\%-level evidence that the presence of an identified article is associated with a different abnormal return. The result is compatible with a modest association, but it is also compatible with sampling variability, omitted market-wide events, or selection into news coverage.

\subsection{Sentiment-Group Contrasts}
\label{subsec:group-tests}

The next analysis compares mean abnormal returns across sentiment categories at the ticker--day level. The three pre-specified contrasts are positive versus neutral, negative versus neutral, and positive versus negative. Each contrast is evaluated with a Welch test, a permutation test, and a bootstrap interval.

\begin{table}[htbp]
\centering
\scriptsize
\caption{Pairwise abnormal-return contrasts across predicted sentiment groups. Differences are computed as group A minus group B.}
\label{tab:sentiment-contrasts}
\begin{tabular}{llrrrrrrrr}
\toprule
\textbf{Contrast} & \textbf{$N_A$} & \textbf{$N_B$} & \textbf{Mean A} & \textbf{Mean B} & \textbf{Difference} & \textbf{Welch $p$} & \textbf{Perm. $p$} & \textbf{Bootstrap 95\% CI} \\
\midrule
Positive $-$ Neutral & 15 & 127 & 0.9068 & 0.0402 & 0.8666 & 0.3253 & 0.3177 & $[-0.7168,2.5064]$ \\
Negative $-$ Neutral & 10 & 127 & -0.3665 & 0.0402 & -0.4068 & 0.7123 & 0.7079 & $[-2.4133,1.5643]$ \\
Positive $-$ Negative & 15 & 10 & 0.9068 & -0.3665 & 1.2733 & 0.3444 & 0.3337 & $[-1.0868,3.7216]$ \\
\bottomrule
\end{tabular}
\end{table}

\begin{figure}[htbp]
\centering
\includegraphics[width=0.86\textwidth]{images/sentiment_group_tests.png}
\caption{Pairwise differences in mean abnormal returns between sentiment groups with bootstrap 95\% confidence intervals. All intervals include zero, and the positive and negative groups are small relative to the neutral group.}
\label{fig:sentiment-group-tests}
\end{figure}

The point estimates have the expected sign only for the same-day positive-minus-neutral comparison and the negative-minus-neutral comparison. However, none of the three contrasts is statistically significant, and every bootstrap interval crosses zero. The results therefore do not support a claim that the predicted sentiment groups have reliably different abnormal returns in this sample. The imbalance in the prediction distribution is a central limitation: the negative group has only ten ticker--day observations and the positive group has fifteen.

\subsection{Two-Way Fixed-Effects Panel Regression}
\label{subsec:panel-regression}

To control for persistent differences between securities and common shocks across dates, we estimate the following two-way fixed-effects specification:
\[
 AR_{i,\tau(i,t)}=\alpha_i+\delta_{\tau(i,t)}+\beta_1S_{i,t}
 +\beta_2\text{negative\_share}_{i,t}
 +\beta_3\text{positive\_share}_{i,t}
 +\gamma\log(1+\text{article\_count}_{i,t})+\varepsilon_{i,t},
\]
where $\tau(i,t)$ is the first trading day strictly after the article date, $\alpha_i$ are ticker fixed effects, and $\delta_{\tau(i,t)}$ are date fixed effects. The dependent variable is market-model abnormal return on the aligned evaluation date. The regression uses 152 ticker--day observations, 51 tickers, and 21 dates.

The intended primary covariance estimator is HC3. In this sparse two-way fixed-effects design, HC3 produces infinite standard errors for some observations because the leverage correction reaches a singular boundary. The pipeline therefore retains the same regression specification and uses HC1 as a documented finite-standard-error fallback. This fallback is reported explicitly rather than silently presenting invalid infinite standard errors. The coefficients should consequently be treated as exploratory conditional associations, not causal estimates.

\begin{table}[htbp]
\centering
\caption{Two-way fixed-effects panel regression for abnormal returns. Standard errors are HC1 because the intended HC3 correction is singular for this sparse design.}
\label{tab:panel-regression}
\begin{tabular}{lrrrr}
\toprule
\textbf{Variable} & \textbf{Coefficient} & \textbf{Std. error} & \textbf{$t$-statistic} & \textbf{$p$-value} \\
\midrule
Intercept & -0.2811 & 1.5913 & -0.1766 & 0.8598 \\
Sentiment score & 5.9079 & 5.3257 & 1.1093 & 0.2673 \\
Negative share & 5.5027 & 5.2564 & 1.0468 & 0.2952 \\
Positive share & -5.6140 & 5.4776 & -1.0249 & 0.3054 \\
Log article count & 0.9852 & 1.3881 & 0.7097 & 0.4779 \\
\midrule
Ticker fixed effects & \multicolumn{4}{l}{Yes} \\
Date fixed effects & \multicolumn{4}{l}{Yes} \\
Observations & \multicolumn{4}{l}{152} \\
Unique tickers & \multicolumn{4}{l}{51} \\
Unique dates & \multicolumn{4}{l}{21} \\
$R^2$ / adjusted $R^2$ & \multicolumn{4}{l}{0.6134 / 0.2419} \\
\bottomrule
\end{tabular}
\end{table}

The sentiment-score coefficient is positive, but its standard error is large relative to the point estimate and its $p$-value is 0.2673. The coefficient therefore does not provide evidence of a statistically significant conditional association after controlling for ticker and date fixed effects. The positive and negative share coefficients should not be interpreted independently as causal marginal effects because shares are compositional measures and are mechanically related to the omitted neutral share. More importantly, the high $R^2$ should not be mistaken for strong predictive evidence: much of the fitted variation can be absorbed by ticker and date fixed effects, while the sentiment covariates remain statistically imprecise.

\subsection{Lagged Robustness and Sensitivity Analysis}
\label{subsec:lag-robustness}

The lag analysis examines whether the abnormal-return pattern persists when the evaluation date is shifted by 0, 1, 2, 3, 5, and 10 trading sessions relative to the aligned event date. These estimates are sensitivity checks rather than independent causal tests. The sample size declines with the lag because later evaluation dates require additional price observations.

\begin{table}[htbp]
\centering
\caption{Lagged abnormal-return robustness. Mean abnormal returns are reported overall and by predicted sentiment.}
\label{tab:lag-robustness}
\begin{tabular}{lrrrrrr}
\toprule
\textbf{Design} & \textbf{Lag} & \textbf{N} & \textbf{Overall} & \textbf{Negative} & \textbf{Neutral} & \textbf{Positive} \\
\midrule
Lag 0 & 0 & 200 & -0.0918 & -0.7884 & -0.1617 & 0.8940 \\
Lag 1 & 1 & 176 & -0.1821 & 0.3368 & 0.0629 & -2.7331 \\
Lag 2 & 2 & 165 & -0.1729 & -1.2344 & -0.0685 & -0.4067 \\
Lag 3 & 3 & 155 & -0.5969 & -2.6554 & -0.4863 & -0.1760 \\
Lag 5 & 5 & 136 & -0.6588 & 0.4479 & -0.8585 & 0.5251 \\
Lag 10 & 10 & 95 & -0.4426 & -0.4422 & -0.2016 & -4.1530 \\
\bottomrule
\end{tabular}
\end{table}

\begin{figure}[htbp]
\centering
\includegraphics[width=0.82\textwidth]{images/robustness_lags.png}
\caption{Lagged sensitivity of mean abnormal returns. The overall estimate varies in magnitude and does not display a stable monotonic decay. Sentiment-specific means are especially unstable because the positive and negative groups are small.}
\label{fig:lag-robustness}
\end{figure}

The overall mean abnormal return is close to zero at lag 0 and becomes more negative at lags 3 and 5. The sign and magnitude of the sentiment-specific means also change across lags: for example, the positive group is positive at lag 0, strongly negative at lag 1, negative at lag 2, and negative again at lag 10. This instability weakens the interpretation of a persistent sentiment-driven price response. The lag results are best reported as evidence that the downstream association is sensitive to the timing convention and available evaluation horizon.

\subsection{Additional Diagnostic Figures}
\label{subsec:diagnostic-figures}

The following figures provide diagnostics for market-model coverage and the relationship between daily sentiment and abnormal return. They are descriptive and should not be read as causal identification.

\begin{figure}[htbp]
\centering
\includegraphics[width=0.82\textwidth]{images/market_model_coverage.png}
\caption{Distribution of the number of pre-event observations used in the market-model estimation. Coverage is heterogeneous across symbols, although most symbols have the full available pre-event history.}
\label{fig:market-model-coverage}
\end{figure}

\begin{figure}[htbp]
\centering
\includegraphics[width=0.92\textwidth]{images/sentiment_vs_abnormal_return.png}
\caption{Daily mean sentiment and mean abnormal return in the ticker--day panel. The figure is a descriptive time-series comparison; it does not adjust for all confounding shocks and does not establish temporal causality.}
\label{fig:sentiment-vs-ar}
\end{figure}

\subsection{Experiments Not Estimated in the Current Run}
\label{subsec:unestimated-experiments}

The original experimental plan also proposed intrinsic manual evaluation, out-of-time predictive benchmarking, probability calibration, and a qualitative error-analysis set. These experiments are not assigned numerical results in this section because the submitted artifacts do not contain the required gold labels, calibrated probabilities, or a separately trained return-prediction benchmark. In particular, the 208 CafeF predictions are model outputs, not human-annotated ground truth; using them as their own labels would create circular evaluation. Similarly, the prediction file contains blank probability columns, so Expected Calibration Error cannot be computed defensibly from the current run.

\begin{table}[htbp]
\centering
\caption{Planned analyses that require additional artifacts before valid estimation.}
\label{tab:deferred-experiments}
\begin{tabular}{p{0.28\textwidth}p{0.58\textwidth}}
\toprule
\textbf{Analysis} & \textbf{Required artifact and reason it is not reported here} \\
\midrule
Intrinsic classification evaluation & Human-annotated labels for a held-out article--ticker sample, preferably with independent annotations and adjudication. Model predictions alone cannot establish accuracy, F1, MCC, or Cohen's kappa. \\
Out-of-time return benchmark & A pre-specified train/validation/test design and fitted predictive models for the return target. The current pipeline is an event-study association analysis, not a trained return forecaster. \\
Probability calibration & Per-class probabilities or logits for every prediction. The current prediction file has empty probability columns. \\
Qualitative error analysis & A verified set of model errors with gold labels and a documented sampling rule. Illustrative examples without gold labels would be anecdotal rather than an evaluation result. \\
Broad-linking robustness & A second full inference and analysis run using the 512-row broad mapping. The current reported results use the 208-row strict mapping only. \\
\bottomrule
\end{tabular}
\end{table}

This separation between completed and deferred experiments is methodologically important. It prevents attractive but unsupported numbers from being presented as empirical evidence and makes clear which additional data would be needed to extend the study.

\subsection{Overall Assessment}
\label{subsec:overall-assessment}

The completed pipeline is internally consistent at the prediction-contract level: all 208 strict inputs have exactly one aligned prediction, all labels are parseable, and article--ticker observations are aggregated to ticker--day units before market outcomes are compared. This is a substantive strength because it prevents row-order joins, duplicate market outcomes, and untracked unknown labels from silently entering the statistical analysis.

The financial results, however, are not statistically conclusive. The same-day CAR for the full sample is close to zero and is not significant under either the one-sample $t$-test or the sign-flip permutation test. The news-day control produces a positive point estimate, but its confidence interval includes zero. Sentiment-group contrasts are imprecise, and the two-way fixed-effects coefficient on the sentiment score is not significant. Lagged estimates vary substantially across horizons and do not show a stable pattern.

The most visible descriptive pattern is that positive predictions have a higher same-day mean abnormal return than neutral predictions, while negative predictions have a lower mean than neutral predictions. This pattern is directionally appealing, but it is not sufficient for a research conclusion because the positive and negative groups contain only 15 and 10 ticker--day observations, respectively, and all relevant confidence intervals include zero. The correct interpretation is therefore that the sample contains a suggestive but unstable association, not evidence that FinABSA sentiment causes or reliably predicts stock-price movements.

\subsection{Conclusion and Limitations}
\label{subsec:validation-conclusion}

For the October 2022 CafeF sample, the analysis does not reject the null hypothesis of zero abnormal return in the aggregate event windows, does not establish a statistically significant difference between news and eligible no-news days, and does not identify a statistically significant sentiment coefficient in the two-way fixed-effects panel. The downstream evidence is consequently negative in the inferential sense: the pipeline successfully produces and tests a sentiment signal, but the available sample does not provide robust statistical support for a stable sentiment--return relationship.

Several limitations qualify this conclusion. First, the sample covers one month and a crisis-like market regime, so the results may be period-specific. Second, CafeF is a single news source, and publication selection is not random. Third, automated entity linking can create both false positives and false negatives; the strict sample improves precision at the cost of coverage. Fourth, the prediction distribution is highly imbalanced toward neutral, leaving little power for the positive and negative comparisons. Fifth, daily prices do not reveal the exact time at which information reached the market, which motivates the next-trading-day alignment but does not eliminate timing error. Sixth, the sparse ticker--day panel makes robust covariance estimation difficult, as illustrated by the HC3-to-HC1 fallback. Finally, the analysis is observational: no coefficient, CAR, or group contrast should be described as a causal effect.

The contribution of this experiment is therefore best framed as an end-to-end NLP-to-finance validation exercise. It demonstrates how target-masked ABSA predictions can be linked to Vietnamese listed securities, transformed into reproducible ticker--day features, and evaluated through multiple complementary statistical tests. It also demonstrates that a technically complete pipeline can yield a scientifically cautious conclusion: a directional descriptive pattern is not equivalent to statistically reliable predictive power.

\begin{table}[htbp]
\centering
\caption{Reproducibility checklist for the completed validation run.}
\label{tab:reproducibility}
\begin{tabular}{p{0.48\textwidth}p{0.38\textwidth}}
\toprule
\textbf{Component} & \textbf{Completed implementation} \\
\midrule
Prediction alignment & Join by stable \texttt{sample\_id}; 208/208 matched \\
Label normalization & Canonical positive/neutral/negative; zero unknown labels \\
Entity-linking primary sample & Strict confidence threshold $\geq 0.95$ \\
Market model & Security-specific market model with VNINDEX proxy \\
Event windows & $[0,0]$, $[0,+1]$, and $[0,+3]$ \\
Inference tests & One-sample $t$-test, sign-flip permutation, bootstrap CI \\
Control test & News-day versus eligible no-news-day comparison \\
Group tests & Welch, permutation, and bootstrap contrasts \\
Panel model & Ticker and date fixed effects; HC1 fallback documented \\
Lag sensitivity & 0, 1, 2, 3, 5, and 10 trading-session lags \\
Generated artifacts & CSV tables, eight figures, Markdown report, PDF report \\
\bottomrule
\end{tabular}
\end{table}

The complete numerical outputs underlying this section are stored in the repository under \texttt{data/cafef\_oct2022/analysis/}. The Markdown and PDF versions of the automatically generated report should be treated as supplementary machine-generated records; the tables and interpretations in this section are the human-reviewed academic presentation of the same run.
